\documentclass[journal]{IEEEtran}
\usepackage{amsmath,amsfonts,amssymb,amsthm}
\usepackage{graphicx}
\usepackage{booktabs}
\usepackage{cite}
\usepackage{algorithm}
\usepackage{algorithmic}
\usepackage{siunitx}
\usepackage{microtype}
\usepackage{subcaption}

\newtheorem{theorem}{Theorem}
\newtheorem{lemma}{Lemma}
\newtheorem{corollary}{Corollary}
\theoremstyle{definition}
\newtheorem{definition}{Definition}

\begin{document}

\title{CCE-QOS: Constraint-Coupled Energy Minimization and Adaptive Penalty Refinement for Combinatorial Operator Scheduling on Neural Processing Units}

\author{Koduru~Yagnesh~Kumar%
\thanks{Koduru Yagnesh Kumar is an independent researcher (email: yagneshkumarkoduru@gmail.com).}}

\maketitle

\begin{abstract}
The energy efficiency of modern Neural Processing Units (NPUs) is critically bounded by data movement across the on-chip SRAM/DRAM memory hierarchy rather than peak arithmetic multiply-accumulate throughput. While state-of-the-art deep learning compilers (e.g., Apache TVM, XLA, MLIR) optimize individual kernel loop nests, end-to-end global operator scheduling across constrained multi-bank SRAM and shared DRAM channels remains governed by greedy heuristics or decoupled phase-ordered passes. These decoupled passes ignore the non-linear coupling between inter-operator tensor reuse, SRAM bank conflicts, and dynamic voltage/frequency scaling (DVFS).

In this paper, we introduce \textbf{CCE-QOS}, a mathematical optimization framework and compiler that formulates NPU operator scheduling as a Constraint-Coupled Energy (CCE) Quadratic Unconstrained Binary Optimization (QUBO) Hamiltonian. To overcome the severe pathology of penalty tuning in constrained binary optimization—where static penalties either generate invalid schedules or collapse the objective landscape—we introduce \textbf{Adaptive Penalty Refinement (APR)}: a dynamic Lagrangian update law with penalty annealing and a feasibility-preserving polish that empirically reduces hard constraint violations (precedence, memory budget, single-execution) on the tested workloads (measured: 51.61\% to 67.74\% feasibility, the maximum across pipeline methods). A general finite-iteration convergence proof is not established; treat APR as an empirically motivated heuristic (see \texttt{EVIDENCE.md}). 
We construct an end-to-end Python compiler pipeline featuring exact McCormick linearization via Google OR-Tools CP-SAT, a variational Quantum Approximate Optimization Algorithm (QAOA) statevector simulation engine with seeded, reproducible parameter optimization and a depth sweep over $p = 1..3$, and an analytical Pareto exploration engine. On the synthetic benchmark workloads in \texttt{metrics.txt} / \texttt{results\_table.txt} (all current-pipeline outputs, regenerated 2026-09-10), the energy-based formulation with APR reaches the measured maximum feasibility of \textbf{67.74\%} (up from 51.61\% for greedy, +16.13pp) while keeping the lowest latency in the CCE-QUBO table (3479.19 cycles); the best classical search in the cost formulation reaches 4168.69 versus the greedy 5669.65 (26.47\% cost reduction). On the energy formulation APR trades 3.9\% higher QUBO energy (160.43 vs 154.40) for the feasibility gain and a 7.5\% lower cost than the same-arm classical Lookahead (6287.24 vs 6794.39). All results are Python benchmarks on example workloads, not production deployments or hardware measurements; see \texttt{EVIDENCE.md} for exact sources.
\end{abstract}

\begin{IEEEkeywords}
Neural Processing Units (NPUs), Combinatorial Optimization, QUBO Hamiltonian, Quantum Approximate Optimization Algorithm (QAOA), Adaptive Penalty Refinement, Compilers.
\end{IEEEkeywords}

\section{Introduction}
\IEEEPARstart{D}{omain-specific} neural processing units (NPUs) have emerged as the foundational compute engine for physical intelligence, autonomous robotics, and edge computer vision. However, the energy cost of accessing off-chip LPDDR memory ($100\text{--}200$~pJ/byte) exceeds on-chip scratchpad SRAM access ($1\text{--}2$~pJ/byte) by two orders of magnitude~\cite{chen2018,horowitz2014}. Under restricted on-chip memory budgets ($0.5\text{--}4.0$~MB), the static execution order of the neural Directed Acyclic Graph (DAG) directly determines buffer liveness, DRAM eviction cascades, and memory bus contention.

Existing compilation frameworks suffer from three systemic limitations:
\begin{enumerate}
    \item \textbf{Decoupled Phase-Ordering Pathologies:} Compilers separate operator fusion, memory allocation, and topological scheduling into serial passes. A scheduling pass that optimizes for critical-path latency frequently fragments SRAM buffers, forcing massive DRAM spilling in subsequent passes.
    \item \textbf{FLOP-Centric Proxy Metrics:} Heuristic schedulers minimize proxy objectives (e.g., total FLOP count or critical path length), ignoring physical CMOS switching dynamics ($E_{\text{dyn}} = C_{\text{eff}} V^2 f$), DVFS state transition penalties, and concurrent SRAM bank contention.
    \item \textbf{Penalty Multiplier Dilemma in Binary Optimization:} Mapping DAG scheduling onto Quadratic Unconstrained Binary Optimization (QUBO) or Ising formulations requires quadratic penalty terms to enforce hard operational constraints. Static penalty multipliers inevitably fail: insufficient penalties produce physically illegal schedules, whereas excessive penalties overwhelm the gradient, trapping classical or quantum solvers in poor local minima.
\end{enumerate}

To resolve these challenges from first principles, we develop \textbf{CCE-QOS}.

\section{Mathematical Formulation}
\subsection{Workload Graph & Binary Decision Variables}
Let an NPU workload be modeled as a directed acyclic graph $\mathcal{G} = (\mathcal{V}, \mathcal{E})$, where vertices $v_i \in \mathcal{V}$ represent tensor operators (e.g., Conv2D, MatMul, LayerNorm) and directed edges $(v_i, v_j) \in \mathcal{E}$ represent tensor dependencies with volume $B(v_i, v_j)$ bytes. The execution timeline is discretized into $T$ sequential slots, and hardware execution modes are denoted by $r \in \mathcal{R}$ (encoding DVFS voltage/frequency pairs).

We define binary decision variables:
\begin{equation}
x_{i,t,r} \in \{0, 1\}, \quad \forall v_i \in \mathcal{V}, \; t \in \{1, \dots, T\}, \; r \in \mathcal{R}
\end{equation}
where $x_{i,t,r} = 1$ if and only if operator $v_i$ commences execution at time slot $t$ under resource mode $r$.

\subsection{Constraint-Coupled Energy (CCE) Hamiltonian}
The total Hamiltonian is formulated as:
\begin{equation}
H_{\text{total}} = H_{\text{unary}} + H_{\text{reuse}} + H_{\text{contention}} + H_{\text{constraints}}
\end{equation}

\subsubsection{Unary Dynamic Energy & DVFS Switching}
\begin{equation}
H_{\text{unary}} = \sum_{i \in \mathcal{V}} \sum_{t=1}^T \sum_{r \in \mathcal{R}} \left( E_{\text{comp}}(v_i, r) + \alpha_{\text{lat}} \tau(v_i, r) \right) x_{i,t,r}
\end{equation}
where $E_{\text{comp}}(v_i, r) = C_{\text{eff}} V_r^2 f_r \cdot \tau(v_i, r)$ models dynamic CMOS switching energy, and $\tau(v_i, r)$ is the execution duration.

\subsubsection{Quadratic Inter-Operator Tensor Reuse}
When a consumer operator $v_j$ executes within the SRAM residency window $\Delta t_{\text{res}}$ of producer $v_i$, intermediate DRAM spill traffic is avoided:
\begin{equation}
H_{\text{reuse}} = -\sum_{(v_i, v_j) \in \mathcal{E}} \sum_{t=1}^T \sum_{\delta=1}^{\Delta t_{\text{res}}} \gamma_{\text{reuse}} \cdot B(v_i, v_j) \cdot \left( \sum_r x_{i,t,r} \right) \left( \sum_{r'} x_{j,t+\delta,r'} \right)
\end{equation}

\subsubsection{SRAM Multi-Bank Contention Penalty}
Concurrent memory accesses to the same physical SRAM bank $k \in \mathcal{K}$ generate quadratic pipeline stalls:
\begin{equation}
H_{\text{contention}} = \sum_{t=1}^T \sum_{k \in \mathcal{K}} \lambda_{\text{bank}} \left( \sum_{i: \text{bank}(v_i)=k} \sum_r x_{i,t,r} \right)^2
\end{equation}

\subsection{Adaptive Penalty Refinement (APR) Theory}
The hard physical constraints encompass:
1) \textbf{Unique Execution:} Each operator must execute exactly once:
\begin{equation}
H_{\text{exec}} = \sum_{i \in \mathcal{V}} \left( \sum_{t=1}^T \sum_{r \in \mathcal{R}} x_{i,t,r} - 1 \right)^2
\end{equation}
2) \textbf{Causal Precedence:} If $(v_i, v_j) \in \mathcal{E}$, $t_i + \tau_i \le t_j$:
\begin{equation}
H_{\text{prec}} = \sum_{(v_i, v_j) \in \mathcal{E}} \sum_{t_i \ge t_j} \left(\sum_r x_{i,t_i,r}\right) \left(\sum_{r'} x_{j,t_j,r'}\right)
\end{equation}
3) \textbf{SRAM Capacity Bound:} Peak active buffer volume must not exceed $M_{\text{SRAM}}$:
\begin{equation}
H_{\text{sram}} = \sum_{t=1}^T \max\left(0, \sum_{i \in \text{Live}(t)} \text{Size}(v_i) - M_{\text{SRAM}}\right)^2
\end{equation}

Under APR, penalty weights $\boldsymbol{\lambda}^{(k)} = [\lambda_{\text{exec}}, \lambda_{\text{prec}}, \lambda_{\text{sram}}]^T$ are updated iteratively:
\begin{equation}
\lambda_m^{(k+1)} = \lambda_m^{(k)} \cdot \left( 1 + \eta_m \cdot \frac{\text{Violations}_m(\mathbf{x}^{(k)})}{\text{Constraints}_m} \right)
\label{eq:apr_law}
\end{equation}

\begin{theorem}[Feasibility threshold for exact minimizers (APR scope)]
Let the feasible set $\mathcal{X}_{\text{feasible}} \subseteq \{0,1\}^N$ be nonempty, let violations be counted in whole units so that $\text{Violations}_m(\mathbf{x}) \ge 1$ whenever constraint $m$ is violated, and let $H_{\text{obj}}(\mathbf{x}) \ge 0$ on the finite search space. Under the APR update law (\ref{eq:apr_law}), if $\hat{\mathbf{x}}(\boldsymbol{\lambda}) = \arg\min_{\mathbf{x}} H(\mathbf{x}; \boldsymbol{\lambda})$ is an exact global minimizer and $\lambda_m > H_{\text{obj}}(\mathbf{x}^*)$ for every $m$, where $\mathbf{x}^* \in \arg\min_{\mathbf{x} \in \mathcal{X}_{\text{feasible}}} H_{\text{obj}}(\mathbf{x})$, then $\hat{\mathbf{x}}(\boldsymbol{\lambda})$ is feasible.
\end{theorem}

\begin{proof}
The feasible schedule $\mathbf{x}^*$ incurs zero penalty, so $H(\mathbf{x}^*; \boldsymbol{\lambda}) = H_{\text{obj}}(\mathbf{x}^*)$. Any infeasible $\mathbf{x}$ violates at least one constraint, so $\text{Violations}_m(\mathbf{x}) \ge 1$ for that $m$ and $H(\mathbf{x}; \boldsymbol{\lambda}) \ge H_{\text{obj}}(\mathbf{x}) + \lambda_m \ge \lambda_m > H_{\text{obj}}(\mathbf{x}^*) = H(\mathbf{x}^*; \boldsymbol{\lambda})$. Hence no infeasible point is a global minimizer.
\end{proof}

Because the law (\ref{eq:apr_law}) multiplies each violated constraint's multiplier by at least $(1 + \eta_m / \text{Constraints}_m)$ per round, the multipliers exceed the threshold $H_{\text{obj}}(\mathbf{x}^*)$ in finitely many rounds whenever violations persist, so an exact inner solver returns a feasible schedule under the stated assumptions. The deployed pipeline, however, uses approximate solvers (CP-SAT with time limits, simulated annealing, QAOA sampling, local-search fallback) for which the exact-minimizer assumption does not hold; no convergence-rate claim is made for the implementation. APR is therefore presented as an empirically motivated heuristic with a feasibility-threshold guarantee in the exact-solver limit, measured at 51.61\% to 67.74\% feasibility on the energy formulation (see \texttt{EVIDENCE.md}). The earlier draft's monotonic finite-iteration convergence claim is withdrawn.

\section{Exact \& Quantum Solvers}
\subsection{McCormick Linearization for Google OR-Tools CP-SAT}
\begin{theorem}[Exact single-term McCormick envelope]
For one binary product term $y_{ij} = x_i x_j$ with $x_i, x_j \in \{0, 1\}$, the continuous relaxation $y_{ij} \le x_i$, $y_{ij} \le x_j$, $y_{ij} \ge x_i + x_j - 1$, $y_{ij} \ge 0$ is the convex hull of the four binary points $(0,0,0)$, $(0,1,0)$, $(1,0,0)$, $(1,1,1)$; its vertices are exactly those four points. This is the classical McCormick construction, cited as textbook background. Exactness holds per product term; for Hamiltonians with many coupled product terms the per-term envelope does not describe the joint relaxation, and the resulting integrality gap is measured empirically rather than assumed zero.
\end{theorem}

We implement this in [`ortools_solver.py`](../../ortools_solver.py); on the tested runner workload CP-SAT returns status OPTIMAL in about $20$~ms (see \texttt{EVIDENCE.md}).

\subsection{Variational QAOA Statevector Engine}
We map the QUBO to an Ising spin Hamiltonian via $x_i = (1 - s_i)/2$ and execute $p$-layer variational evolution:
\begin{equation}
|\psi(\boldsymbol{\gamma}, \boldsymbol{\beta})\rangle = \prod_{l=1}^p e^{-i \beta_l \sum_j \sigma_x^{(j)}} e^{-i \gamma_l H_C} |+\rangle^{\otimes n}
\end{equation}
The engine optimizes variational parameters $(\boldsymbol{\gamma}, \boldsymbol{\beta})$ via numerical finite-difference gradients in pure NumPy ([`QAOA_solver.py`](../../QAOA_solver.py)) and exports OpenQASM 2.0 quantum circuits. Parameter initialization is seeded (\texttt{DEFAULT\_QAOA\_SEED} $= 12345$) so repeated optimizations return bit-identical $(\boldsymbol{\gamma}, \boldsymbol{\beta}, E)$ tuples, and the engine supports any depth $p \ge 1$ (parameter vector length $2p$) with a COBYLA depth sweep (\texttt{sweep\_qaoa\_depth}). On the solver's reference 3-qubit instance, the measured ground-state approximation ratio at $p=2$ is $r = 0.8012$ (deterministic, seeded run; replaces the stochastic $r = 0.611$ figure), and the $p = 1..3$ sweep (COBYLA, 60 evaluations per depth) improves it to $r = 0.9385$ at $p=3$. On a real 11-variable workload sub-instance (the \texttt{conv\_stem} $\to$ \texttt{bn\_stem} chain from \texttt{example\_workload.json}), the sweep improves the ratio from $r = 0.7244$ ($p=1$; the old default $p=1$ gradient run measures $r = 0.7088$) to $r = 0.7823$ at $p=3$.

\subsection{LLM KV-Cache Paging \& Continuous Batching Extension}
To support modern Large Language Model (LLM) inference on edge NPUs, CCE-QOS includes a deterministic Key-Value (KV) cache capacity simulation. The simulator in \texttt{implementations/v3\_llm\_kvcache\_continuous\_batching/llm\_kvcache\_paging\_scheduler.py} compares static maximum-length reservation with exact-demand paged admission on a 50-request synthetic trace and a 1600-block pool of 16-token blocks. Static reservation records 17.4031 output tokens per simulated step and P95 end-to-end latency 402.85 steps; exact-demand admission records 30.2720 output tokens per step and P95 latency 242.20 steps, a 1.7395$\times$ modeled throughput ratio and 39.878\% lower P95 latency. The model queues work safely when capacity is exhausted and does not simulate vLLM behavior, GPU or NPU hardware, host paging, or DRAM faults. A separate local vLLM attribution benchmark is scoped in \texttt{EVIDENCE.md}; on its recorded synthetic trace, stock prefix caching is the dominant effect and the bounded dispatcher does not improve over FIFO plus stock prefix caching.

\section{Experimental Results}
\subsection{Comparative Energy \& Scheduling Performance}
We evaluate CCE-QOS on synthetic example workloads (\texttt{example\_workload.json}, \texttt{outputs/schedules.json}) against greedy, simulated annealing, lookahead, and statevector-QAOA baselines in Table~\ref{tab:compiler}. No TVM, XLA, or MLIR comparison exists in this draft.

\begin{table}[ht]
\centering
\caption{End-to-End NPU Compiler Scheduling Comparison (current pipeline, 2026-09-10)}
\label{tab:compiler}
\begin{tabular}{lcccc}
\toprule
\textbf{Compiler / Strategy} & \textbf{Cost} & \textbf{Energy (QUBO)} & \textbf{Latency (cycles)} & \textbf{Feasible (\%)} \\
\midrule
Greedy Topological (cost formulation) & 5669.65 & 5669.65 & 3456.25 & 51.61 \\
Simulated Annealing (cost formulation) & 4443.10 & 4443.10 & 3488.46 & 51.61 \\
\textbf{Lookahead Tree Search (best classical, cost formulation)} & \textbf{4168.69} & 4168.69 & 3390.81 & 54.83 \\
Lookahead / Beam Search (best classical on the CCE-QUBO energy) & 6794.39 & \textbf{154.40} & 3603.50 & 51.61 \\
Simulated Annealing (CCE-QUBO energy) & 6461.85 & 159.31 & 3563.48 & 58.06 \\
\textbf{CCE + APR (energy formulation, ours)} & 6287.24 & 160.43 & \textbf{3479.19} & \textbf{67.74} \\
Quantum (local-search fallback, honestly labeled) & 9756.86 & 164.94 & 3834.24 & 45.16 \\
\bottomrule
\end{tabular}
\end{table}

Measured, current pipeline (all figures from \texttt{metrics.txt} / \texttt{results\_table.txt}, regenerated 2026-09-10; root and \texttt{outputs/} copies are identical): the best classical search reaches \textbf{4168.69} versus greedy 5669.65 (\textbf{26.47\% cost reduction}) in the cost formulation, and APR with penalty annealing and a feasibility-preserving polish improves feasibility \textbf{51.61\% to 67.74\%} (+16.13pp, the measured maximum) on the energy formulation, with cost 6287.24 (7.5\% lower than the same-arm Lookahead 6794.39) and the lowest latency in the CCE-QUBO table (3479.19 cycles); APR trades 3.9\% higher QUBO energy (160.43 vs 154.40) for these gains. The example workload's QUBO has 4247 variables, so statevector QAOA is not run there; the QAOA engine itself is measured on the reference 3-qubit instance and an 11-variable workload sub-instance (see the Variational QAOA Statevector Engine section). Feasibility is the measured maximum (67.74\%), not 100\%. Latency values are scheduler-model cycles, not hardware measurements. The earlier 4216.92 / 25.62\% / X = 0.2562 figures came from a retired pipeline and are withdrawn (see \texttt{EVIDENCE.md}).

\subsection{Pareto Frontier Analysis}
The 13-point trade-off frontier and the 22.4\% SRAM residency claim require dominance verification against output files before reuse (see \texttt{EVIDENCE.md}). No figure is embedded in this draft.

\section{Conclusion}
CCE-QOS formulates NPU operator scheduling as a coupled quadratic Hamiltonian with Adaptive Penalty Refinement and evaluates it as a Python benchmark pipeline on synthetic example workloads. Measured results show the best classical search at a 26.47\% cost reduction versus greedy in the cost formulation (4168.69 vs 5669.65) and a 51.61\%-to-67.74\% feasibility improvement on the energy formulation via APR with penalty annealing and a feasibility-preserving polish. Production deployment, SOTA-compiler comparison, hardware measurement, and full feasibility guarantees are explicitly out of scope for this draft; see \texttt{EVIDENCE.md}.

\bibliographystyle{IEEEtran}
\begin{thebibliography}{99}
\bibitem{chen2018} T. Chen et al., ``TVM: An automated end-to-end optimizing compiler for deep learning,'' in \emph{USENIX OSDI}, 2018, pp. 578--594.
\bibitem{horowitz2014} M. Horowitz, ``1.1 Computing's energy problem (and what we can do about it),'' in \emph{IEEE ISSCC}, 2014, pp. 10--14.
\bibitem{farhi2014} E. Farhi et al., ``A quantum approximate optimization algorithm,'' \emph{arXiv:1411.4028}, 2014.
\bibitem{lucas2014} A. Lucas, ``Ising formulations of many NP problems,'' \emph{Frontiers in Physics}, vol. 2, p. 5, 2014.
\end{thebibliography}

\end{document}
